\documentclass[11pt]{article}
\usepackage{amsmath,amsthm,amssymb}
\usepackage[margin=1in]{geometry}
\usepackage{enumitem}
\usepackage{hyperref}

\newtheorem{theorem}{Theorem}
\newtheorem{lemma}[theorem]{Lemma}
\newtheorem{proposition}[theorem]{Proposition}
\newtheorem{corollary}[theorem]{Corollary}
\theoremstyle{definition}
\newtheorem{definition}[theorem]{Definition}
\newtheorem{remark}[theorem]{Remark}

\title{The hook tracer coefficient:\\
       closed form $g_k(q) = (q+1)^{\binom{k}{2}}$}
\author{Clio}
\date{13 May 2026 (after very-late-evening)}

\begin{document}
\maketitle

\begin{abstract}
For every hook $\lambda = (k, 1^{n-k})$ in the rank-zero regime $n \ge 2k$,
the matrix coefficient
\[
   g_k(q) \;:=\; \langle v_{T_{*}^{(\lambda)}},\;
   B^{(\lambda)}_{j_0(\lambda)}\,v_{T_{\mathrm{first}}^{(\lambda)}}\rangle
\]
between the two tracer SYTs of \texttt{2026-05-13-hook-dim-equals-one.tex}
equals
\[
   g_k(q) \;=\; (q+1)^{\binom{k}{2}},
\]
in the asymmetric Murphy normalization of the seminormal basis (the
convention $b' = 1$ for every 2-block off-diagonal).  This upgrades the
``$g_k(q) \ne 0$'' theorem of the hook-dim-1 paper to an exact
closed form, settling its Corollary~5.2 (the explicit-formula
conjecture, verified there only computationally to $k = 5$).

\textbf{The proof refines Lemmas A and B of the hook-dim-1 paper.}  In
the asymmetric Murphy normalization, the recursion
$g_k = C_k \cdot g_{k-1} \cdot f_k$ specialises to
$C_k = (q+1)^{k-1}$ and $f_k = 1$.  The first equality is the
$R'_n$-trace through the same-row prefix (which contributes
$(q+1)^{k-1}$) followed by a chain of 2-blocks in which the off-diagonal
$b' = 1$ propagates the $(q+1)^{k-1}$ coefficient unchanged onto
$v_{T'_-}$.  The second equality is the all-swap chain
$T'_- \to \widetilde{T}^{(1)} \to \cdots \to \widetilde{T}^{(k-2)} = T_{*}^{(\lambda)}$
inside $M v_{T'_-}$, where the off-diagonal $b' = 1$ at each step
contributes the multiplicative factor $1$.  Solving the recursion with
the base case $g_2(q) = q + 1$ gives the stated closed form.

The result depends on the choice of normalization of the seminormal
basis, but its existence as a single $(q+1)$-monomial does not.  In
particular, it implies a clean numerical identity:
$g_k(7) = 8^{\binom{k}{2}}$, recovered by the verified data of the
hook-dim-1 paper.
\end{abstract}

\section{Setup}\label{sec:setup}

\subsection{Seminormal normalization}

We work in the seminormal basis of $V_\lambda$ used throughout the
May-2026 stream of papers (the asymmetric Murphy convention from
\texttt{verify\_seminormal.py}).  On each $T_i$-2-block
$\mathrm{span}(v_T, v_{s_iT})$ at axial distance $d \ne \pm 1$, the
matrix of $T_i$ is
\begin{equation}\label{eq:hoefsmit}
   T_i\,(v_T,\,v_{s_iT}) \;=\;
   (v_T,\,v_{s_iT})\,
   \begin{pmatrix} a_d & b_d b'_d \\ 1 & a'_d \end{pmatrix},
\end{equation}
with diagonal entries
\[
   a_d \;=\; \frac{(q-1)q^d}{q^d - 1}, \qquad
   a'_d \;=\; \frac{q-1}{1 - q^d}\;=\; -\,a_d/q^d,
\]
off-diagonal $b'_d = 1$, and the remaining off-diagonal
$b_d b'_d = a_d a'_d + q = q - a_d^2 / q^d$.  The eigenvalues of $T_i$
on this 2-block are $q$ and $-1$.

In this normalization, the action of $(T_i + 1)$ on the 2-block is:
\begin{align*}
   (T_i + 1) v_T &\;=\; (a_d + 1)\,v_T + 1\cdot v_{s_iT}, \\
   (T_i + 1) v_{s_iT} &\;=\; b_d b'_d \cdot v_T + (a'_d + 1)\,v_{s_iT}.
\end{align*}
A direct computation in $q$ gives the identity
\begin{equation}\label{eq:eigvec}
   q - a'_d \;=\; a_d + 1,
\end{equation}
so the $+q$-eigenvector of $T_i$ on the 2-block, normalized to have
coefficient $1$ on $v_{s_iT}$, is $(a_d + 1)\, v_T + v_{s_iT}$.

\subsection{Tracer SYTs and the conjectural formula}

For $\lambda = (k, 1^{n-k})$, $k \ge 2$, $n \ge 2k$, we use the
hook-$S$ notation: an SYT of $\lambda$ is determined by its
$(k-1)$-element subset $S \subset \{2,\ldots,n\}$ filling row $1$ after
the entry $1$ at $(1,1)$.  The two tracer SYTs of
\texttt{2026-05-13-hook-dim-equals-one.tex} are
\[
   T_{\mathrm{first}}^{(\lambda)}: \; S = \{2, 3, \ldots, k\},
   \qquad
   T_{*}^{(\lambda)}: \; S = \{n - k + 2, \ldots, n\}.
\]
The tracer coefficient is
\[
   g_k(q) \;=\; \langle v_{T_{*}^{(\lambda)}},\;
   B^{(\lambda)}_{j_0(\lambda)}\, v_{T_{\mathrm{first}}^{(\lambda)}}\rangle,
\]
where $B^{(\lambda)}_{j_0} = R'_{j_0}\,R'_{j_0+1}\,\cdots\,R'_n$ with
$j_0(\lambda) = n - k + 2$ and
$R'_k = (T_{k-1}+1)(T_{k-2}+1)\cdots(T_1+1)$.

\begin{theorem}[Main]\label{thm:main}
For every $k \ge 2$ and every $n \ge 2k$,
\[
   g_k(q) \;=\; (q+1)^{\binom{k}{2}}.
\]
In particular $g_k$ is a polynomial in $\mathbb{Z}_{\ge 0}[q]$
independent of $n$ inside the rank-zero regime.
\end{theorem}

\subsection{The recursion}

The hook-dim-1 paper's Lemma A and Lemma B set up a recursion
\begin{equation}\label{eq:recursion-frame}
   g_k(q) \;=\; C_k \cdot g_{k-1}(q) \cdot f_k(q),
\end{equation}
where:
\begin{itemize}[leftmargin=2em,topsep=0.2em,itemsep=0.1em]
\item $C_k$ is the scalar in Lemma~A: $R'_n\,v_{T_{\mathrm{first}}^{(\lambda)}}
   = C_k \cdot \Phi_*(v_{T_{\mathrm{first}}^{(\mu_1)}})$, where
   $\mu_1 = (k-1, 1^{n-k-1})$ and $\Phi_*$ is the embedding from
   Definition~2 of the hook-dim-1 paper;
\item $f_k$ is the scalar in Lemma~B:
   $\langle v_{T_{*}^{(\lambda)}}, M\,\Phi_*(v_{T_{*}^{(\mu_1)}})\rangle = f_k$,
   with $M = (T_\ell{+}1)(T_{\ell+1}{+}1)\cdots(T_{n-2}{+}1)$ and
   $\ell = n - k + 1 = j_0(\lambda) - 1$.
\end{itemize}
We sharpen the qualitative statements ``$C_k \ne 0$, $f_k \ne 0$'' to
the quantitative
\[
   C_k = (q+1)^{k-1}, \qquad f_k = 1
\]
in the normalization fixed above.

\section{Refined Lemma A: $C_k = (q+1)^{k-1}$}\label{sec:lemA}

We normalize $\Phi_*$ by
\begin{equation}\label{eq:phi-norm}
   \Phi_*(v_{T'}) \;=\; (a_{-(n-1)} + 1)\, v_{T'_+} + 1 \cdot v_{T'_-},
\end{equation}
the $+q$-eigenvector of $T_{n-1}$ on the corner pair
$\{(1,k), (n-k+1,1)\}$ (axial distance $-(n-1)$), scaled to have
coefficient $1$ on $v_{T'_-}$.  This is consistent with the convention
of the hook-dim-1 paper (Definition~2) once a $b'$-normalization is
fixed.

\begin{lemma}[Refined Lemma A]\label{lem:Ck}
For every $k \ge 2$ and $n \ge 2k$,
\[
   R'_n\,v_{T_{\mathrm{first}}^{(\lambda)}}
   \;=\; (q+1)^{k-1} \cdot \Phi_*(v_{T_{\mathrm{first}}^{(\mu_1)}}).
\]
In particular, $C_k = (q+1)^{k-1}$.
\end{lemma}

\begin{proof}
We trace $R'_n = (T_{n-1}{+}1)(T_{n-2}{+}1)\cdots(T_1{+}1)$ applied
right-to-left to $v_{T_{\mathrm{first}}^{(\lambda)}}$, identifying the
seminormal support and coefficients at each step.

\textbf{Same-row prefix (steps $1$ through $k-1$).}  For each
$j \in [1, k-1]$, entries $j$ and $j+1$ both lie in row $1$ of
$T_{\mathrm{first}}^{(\lambda)}$ at adjacent cells $(1,j)$ and $(1,j+1)$.
Hence $(T_j + 1) v_{T_{\mathrm{first}}^{(\lambda)}} = (q+1)\, v_{T_{\mathrm{first}}^{(\lambda)}}$.
Cumulatively:
\[
   (T_{k-1}+1)\cdots(T_1+1)\, v_{T_{\mathrm{first}}^{(\lambda)}}
   \;=\; (q+1)^{k-1}\, v_{T_{\mathrm{first}}^{(\lambda)}}.
\]

\textbf{2-block phase (steps $k$ through $n-1$).}  Define, for $m \ge 1$,
the SYT $T^{(m)}$ of $\lambda$ by the row-$1$ data
\[
   T^{(m)}: \; S = \{2, 3, \ldots, k-1, k+m\}.
\]
Thus $T^{(0)} = T_{\mathrm{first}}^{(\lambda)}$, and $T^{(m)}$ is
obtained from $T^{(m-1)}$ by transposing the entries $k+m-1$ and $k+m$
(this is $s_{k+m-1} T^{(m-1)}$ on tableaux).  The positions in $T^{(m)}$
are:
\begin{itemize}[leftmargin=2em,topsep=0.2em,itemsep=0.1em]
\item Row 1: $(1, 1) = 1, (1, 2) = 2, \ldots, (1, k-1) = k-1, (1, k) = k+m$.
\item Column 1: $(2, 1) = k, (3, 1) = k+1, \ldots, (m+1, 1) = k+m-1,
   (m+2, 1) = k+m+1, \ldots, (n-k+1, 1) = n$.
\end{itemize}
We claim by induction on $j$ that for $j = k, k+1, \ldots, n-1$,
\begin{equation}\label{eq:state-j}
   (T_j+1)\cdots(T_1+1)\, v_{T_{\mathrm{first}}^{(\lambda)}}
   \;=\; \alpha_j\, v_{T^{(j-k)}} + (q+1)^{k-1}\, v_{T^{(j-k+1)}},
\end{equation}
for some $\alpha_j \in \mathbb{Q}(q)$ whose precise form we record below.

\textbf{Step $j = k$ (base of induction).}  Apply $(T_k + 1)$ to
$(q+1)^{k-1}\, v_{T_{\mathrm{first}}}$.  In $T_{\mathrm{first}} = T^{(0)}$,
entry $k$ is at $(1, k)$ (content $k-1$) and entry $k+1$ is at $(2, 1)$
(content $-1$).  Axial distance $d_k = -1 - (k-1) = -k \ne \pm 1$, so
$T_k$ acts as a 2-block with partner $T^{(1)}$.  Using~\eqref{eq:hoefsmit}:
\[
   (T_k+1)\, v_{T^{(0)}}
   \;=\; (a_{-k} + 1)\,v_{T^{(0)}} + 1 \cdot v_{T^{(1)}},
\]
hence
\[
   (T_k+1)\cdots(T_1+1)\, v_{T_{\mathrm{first}}^{(\lambda)}}
   \;=\; (q+1)^{k-1}(a_{-k} + 1)\,v_{T^{(0)}} + (q+1)^{k-1}\,v_{T^{(1)}},
\]
matching~\eqref{eq:state-j} with $\alpha_k = (q+1)^{k-1}(a_{-k}+1)$.

\textbf{Inductive step $j \to j+1$, $j \in [k, n-2]$.}  Suppose
\eqref{eq:state-j} holds at step $j$.  Apply $(T_{j+1} + 1)$.  By the
position data of $T^{(m)}$ above:
\begin{itemize}[leftmargin=2em,topsep=0.2em,itemsep=0.1em]
\item In $T^{(j-k)}$, entry $j+1$ is at $(j-k+2, 1)$ and entry $j+2$ is
   at $(j-k+3, 1)$ (both in column~1, adjacent).  Axial distance $-1$:
   $(T_{j+1} + 1) v_{T^{(j-k)}} = 0$.
\item In $T^{(j-k+1)}$, entry $j+1$ is at $(1, k)$ (content $k-1$) and
   entry $j+2$ is at $(j-k+3, 1)$ (content $-(j-k+2)$, from the column-1
   description of $T^{(j-k+1)}$: $(2,1) = k, \ldots, (j-k+2, 1) = j,
   (j-k+3, 1) = j+2$, with the gap at the row-1 cell $(1, k) = j+1$).
   Axial distance
   $d_{j+1} = -(j-k+2) - (k-1) = -(j+1)$.  Hence $T_{j+1}$ acts as a
   2-block with partner $T^{(j-k+2)}$:
   \[
      (T_{j+1}+1)\,v_{T^{(j-k+1)}}
      \;=\; (a_{-(j+1)} + 1)\,v_{T^{(j-k+1)}} + 1 \cdot v_{T^{(j-k+2)}}.
   \]
\end{itemize}
Combining: the state at step $j+1$ equals
\[
   (q+1)^{k-1}(a_{-(j+1)}+1)\,v_{T^{(j-k+1)}} + (q+1)^{k-1}\, v_{T^{(j-k+2)}},
\]
which matches~\eqref{eq:state-j} at $j+1$ with
$\alpha_{j+1} = (q+1)^{k-1}(a_{-(j+1)}+1)$ and the same persistent
coefficient $(q+1)^{k-1}$ on the new SYT.  This completes the induction:
$\alpha_j = (q+1)^{k-1}(a_{-j}+1)$ for every $j \in [k, n-1]$.

\textbf{The final step $j = n-1$.}  At step $n-1$, we apply
$(T_{n-1}+1)$ to the state from step $n-2$.  By~\eqref{eq:state-j} at
$j = n - 2$, the input is supported on $\{T^{(n-k-2)}, T^{(n-k-1)}\}$.
The state after applying $(T_{n-1}+1)$ is supported on
$\{T^{(n-k-1)}, T^{(n-k)}\}$, with coefficients (per the inductive
formula at $j = n-1$):
\[
   R'_n\, v_{T_{\mathrm{first}}^{(\lambda)}}
   \;=\; (q+1)^{k-1}(a_{-(n-1)}+1)\, v_{T^{(n-k-1)}}
       + (q+1)^{k-1}\, v_{T^{(n-k)}}.
\]

\textbf{Identification with $\Phi_*$.}  The tableaux $T^{(n-k-1)}$ and
$T^{(n-k)}$ are precisely $T'_+$ and $T'_-$ for
$T' = T_{\mathrm{first}}^{(\mu_1)}$:
\begin{itemize}[leftmargin=2em,topsep=0.2em,itemsep=0.1em]
\item $T^{(n-k-1)}$: $S = \{2, \ldots, k-1, n-1\}$, so entry $n-1$ at
   $(1,k)$ and entry $n$ at $(n-k+1, 1)$.  This is $T'_+$.
\item $T^{(n-k)}$: $S = \{2, \ldots, k-1, n\}$, so entry $n$ at $(1,k)$
   and entry $n-1$ at $(n-k+1, 1)$.  This is $T'_-$.
\end{itemize}
Therefore
\[
   R'_n\, v_{T_{\mathrm{first}}^{(\lambda)}}
   \;=\; (q+1)^{k-1} \cdot \big[\,(a_{-(n-1)} + 1)\, v_{T'_+}
       + 1 \cdot v_{T'_-}\,\big]
   \;=\; (q+1)^{k-1} \cdot \Phi_*(v_{T_{\mathrm{first}}^{(\mu_1)}}),
\]
using the normalization~\eqref{eq:phi-norm} of $\Phi_*$.
\end{proof}

\section{Refined Lemma B: $f_k = 1$}\label{sec:lemB}

\begin{lemma}[Refined Lemma B]\label{lem:fk}
For every $k \ge 2$ and $n \ge 2k$, with the normalization of
$\Phi_*$ in~\eqref{eq:phi-norm}:
\[
   \langle v_{T_{*}^{(\lambda)}},\; M\,\Phi_*(v_{T_{*}^{(\mu_1)}})\rangle
   \;=\; 1.
\]
In particular, $f_k = 1$.
\end{lemma}

\begin{proof}
By Step 1 of Lemma B in the hook-dim-1 paper, only the $v_{T'_-}$
summand of $\Phi_*(v_{T_{*}^{(\mu_1)}}) = (a_{-(n-1)}+1) v_{T'_+} +
v_{T'_-}$ contributes to the coefficient at $v_{T_{*}^{(\lambda)}}$:
the factors of $M$ do not involve $T_{n-1}$, so entry $n$ is fixed under
$M$; in $T'_+$ it is at $(n-k+1, 1)$, in $T_{*}^{(\lambda)}$ it is at
$(1, k)$.  Hence
\[
   \langle v_{T_{*}^{(\lambda)}},\, M\,\Phi_*(v_{T_{*}^{(\mu_1)}})\rangle
   \;=\; 1 \cdot \langle v_{T_{*}^{(\lambda)}},\, M\, v_{T'_-}\rangle.
\]

\textbf{The $T'_-$ data.}  For $T' = T_{*}^{(\mu_1)}$, the tableau
$T'_-$ has row $1 = \{1, n-k+1, n-k+2, \ldots, n-2, n\}$ and column
$1 = \{1, 2, 3, \ldots, n-k, n-1\}$.  Explicitly,
\[
   (1, j) = n-k+j-1 \text{ for } j \in [2, k-1], \quad (1, k) = n,
\]
\[
   (i, 1) = i \text{ for } i \in [1, n-k], \quad (n-k+1, 1) = n-1.
\]

\textbf{The all-swap chain.}  Define the SYTs
$\widetilde{T}^{(0)} = T'_-$ and inductively for $m = 1, 2, \ldots, k-2$:
\[
   \widetilde{T}^{(m)} \;:=\; s_{n-1-m}\, \widetilde{T}^{(m-1)},
\]
i.e., $\widetilde{T}^{(m)}$ is obtained from $\widetilde{T}^{(m-1)}$ by
swapping the entries $n-1-m$ and $n-m$.  Tracking positions:

\begin{lemma}[Position tracking]\label{lem:positions}
For each $m \in \{0, 1, \ldots, k-2\}$, the tableau $\widetilde{T}^{(m)}$
has:
\begin{itemize}[leftmargin=2em,topsep=0.2em,itemsep=0.1em]
\item Row 1: $(1, j) = n-k+j-1$ for $j \in [2, k-m-1]$,
   $(1, k-i) = n-i$ for $i \in [1, m]$, and $(1, k) = n$.
\item Column 1: $(i, 1) = i$ for $i \in [1, n-k]$,
   $(n-k+1, 1) = n-1-m$.
\end{itemize}
In particular, $\widetilde{T}^{(k-2)}$ has row $1 = \{1, n-k+2,
n-k+3, \ldots, n-1, n\}$ and column $1 = \{1, 2, \ldots, n-k+1\}$,
which is precisely $T_{*}^{(\lambda)}$.
\end{lemma}

\begin{proof}[Proof of Lemma~\ref{lem:positions}]
By induction on $m$.  Base $m = 0$: the description matches $T'_-$
above.  Inductive step: applying $s_{n-1-m}$ to $\widetilde{T}^{(m-1)}$
swaps entries $n-1-m$ and $n-m$.  In $\widetilde{T}^{(m-1)}$, entry
$n-m$ sits at $(n-k+1, 1)$ (by the IH at $m - 1$, since $n - 1 - (m-1)
= n - m$), and entry $n-1-m$ sits at $(1, k - m)$ (the row-$1$ slot
not yet ``rotated up'').  After the swap, $(1, k-m) = n-m$ and
$(n-k+1, 1) = n-1-m$, matching the claim at $m$.  At $m = k-2$ the
row-1 entries are exactly $\{1, n-k+2, n-k+1+1, \ldots, n-1, n\} =
\{1, n-k+2, n-k+3, \ldots, n-1, n\}$ (in order), and the column-1
bottom is $n-k+1$, which is $T_{*}^{(\lambda)}$.
\end{proof}

\textbf{$M$ acts on the all-swap chain by off-diagonal $b' = 1$.}  Write
$M = (T_\ell{+}1)(T_{\ell+1}{+}1)\cdots(T_{n-2}{+}1)$ applied right-to-left
(i.e., $(T_{n-2}{+}1)$ first, $(T_\ell{+}1)$ last).  Index steps of $M$
by $m = 1, 2, \ldots, k-2$, with step $m$ applying the Hecke generator
$T_{n-1-m}$.  By Lemma~\ref{lem:positions} applied to
$\widetilde{T}^{(m-1)}$, the entries $n-1-m$ and $n-m$ sit at row $1$
cell $(1, k-m)$ and column $1$ bottom $(n-k+1, 1)$ respectively, with
axial distance
\[
   d_m \;=\; c(n-m) - c(n-1-m)
   \;=\; -(n-k) - (k-m-1) \;=\; -(n-1-m).
\]
Since $n \ge 2k$ and $m \le k - 2$, we have $n - 1 - m \ge 2$, so
$d_m \le -2$ — a non-degenerate 2-block.  Hence
\begin{equation}\label{eq:swap-step}
   (T_{n-1-m} + 1)\, v_{\widetilde{T}^{(m-1)}}
   \;=\; (a_{d_m} + 1)\, v_{\widetilde{T}^{(m-1)}}
       + 1 \cdot v_{\widetilde{T}^{(m)}}.
\end{equation}

\textbf{Support and uniqueness of the all-swap branch.}  We claim that
for each $m \in \{0, 1, \ldots, k-2\}$, the support of
$(T_{n-1-m}{+}1)\cdots(T_{n-2}{+}1)\, v_{T'_-}$ is exactly
$\{v_{\widetilde{T}^{(0)}}, v_{\widetilde{T}^{(1)}}, \ldots,
v_{\widetilde{T}^{(m)}}\}$.

Indeed, by induction: at step $m$, the input support is
$\{\widetilde{T}^{(0)}, \ldots, \widetilde{T}^{(m-1)}\}$.  We split
into cases by $j$.

\emph{Case $j = m - 1$} (the newest SYT, on which $T_{n-1-m}$ has been
computed as a 2-block by~\eqref{eq:swap-step}): the 2-block produces
$\widetilde{T}^{(m)}$ as the partner with off-diagonal coefficient
$b' = 1$ (the new SYT).

\emph{Case $j \le m - 2$} (older SYTs): we claim the action is
same-row.  Apply Lemma~\ref{lem:positions} to $\widetilde{T}^{(j)}$.
Since $m \ge j + 2$, both $n-1-m$ and $n-m$ are strictly less than the
down-shifted entry $n-1-j$ at $(n-k+1, 1)$, and both are strictly less
than the touched-range minimum $n-j$ (which is $\ge n - m + 1$ here).
Hence both entries lie in the untouched row-$1$ range $(1, j') = n-k+j'-1$
for $j' \in [2, k-j-1]$.  Solving:
\begin{align*}
   n-1-m &= n-k+j'-1 \;\Rightarrow\; j' = k - m, \\
   n-m &= n-k+j'-1 \;\Rightarrow\; j' = k - m + 1.
\end{align*}
Both $j' = k-m$ and $k-m+1$ lie in $[2, k-j-1]$ since $m \in [j+2, k-2]$.
So in $\widetilde{T}^{(j)}$, entry $n-1-m$ is at $(1, k-m)$ and entry
$n-m$ is at $(1, k-m+1)$.  These are adjacent row-1 cells, hence
$(T_{n-1-m} + 1)\, v_{\widetilde{T}^{(j)}} = (q+1)\,
v_{\widetilde{T}^{(j)}}$.  No new SYT introduced.

Hence after step $m$, the support is
$\{\widetilde{T}^{(0)}, \ldots, \widetilde{T}^{(m)}\}$, as claimed.
In particular, $\widetilde{T}^{(m)}$ is the \emph{unique} new SYT, and
its coefficient picks up only from the off-diagonal
$b' = 1$ in~\eqref{eq:swap-step} applied to $\widetilde{T}^{(m-1)}$.

\textbf{Computing the coefficient at $T_{*}^{(\lambda)} = \widetilde{T}^{(k-2)}$.}
By the induction, $v_{\widetilde{T}^{(k-2)}}$ first appears at step
$k - 2$ with coefficient $1$ (the off-diagonal $b'$ of~\eqref{eq:swap-step})
times the coefficient of $\widetilde{T}^{(k-3)}$ at step $k - 3$,
which is in turn $1$ times the coefficient of $\widetilde{T}^{(k-4)}$
at step $k - 4$, etc.  Since each transition picks up exactly $b' = 1$
and the initial coefficient of $\widetilde{T}^{(0)} = T'_-$ is $1$,
we get
\[
   \langle v_{T_{*}^{(\lambda)}},\; M\,v_{T'_-}\rangle
   \;=\; \prod_{m=1}^{k-2} 1 \;=\; 1.
\]
(For $k = 2$, $M$ is the empty product, so $M v_{T'_-} = v_{T'_-}$ and
$T'_- = T_{*}^{(\lambda)}$ directly, giving coefficient $1$ as well.)
\end{proof}

\section{Proof of Theorem~\ref{thm:main}}\label{sec:main}

\begin{proof}[Proof of Theorem~\ref{thm:main}]
Induction on $k$.

\textbf{Base case $k = 2$.}  $\lambda = (2, 1^{n-2})$, $j_0(\lambda) = n$.
$T_{\mathrm{first}}^{(\lambda)}$ has $S = \{2\}$; $T_{*}^{(\lambda)}$ has
$S = \{n\}$.  Apply Lemma~\ref{lem:Ck} with $k = 2$: $R'_n\,
v_{T_{\mathrm{first}}^{(\lambda)}} = (q+1)\cdot \Phi_*(v_{T_{\mathrm{first}}^{(\mu_1)}})$.
But for $k = 2$, $\mu_1 = (1, 1^{n-3}) = (1^{n-2})$ is the
$(n-2)$-sign representation; its unique SYT is
$T_{\mathrm{first}}^{(\mu_1)} = T_{*}^{(\mu_1)}$, and
$\Phi_*(v_{T'}) = (a_{-(n-1)}+1) v_{T'_+} + v_{T'_-}$ with
$T'_- = T_{*}^{(\lambda)}$.  Hence
\[
   g_2(q) \;=\; \langle v_{T_{*}^{(\lambda)}},\, R'_n\,v_{T_{\mathrm{first}}^{(\lambda)}}\rangle
   \;=\; (q+1) \cdot 1 \;=\; q+1.
\]
This matches $(q+1)^{\binom{2}{2}} = (q+1)^1 = q+1$.

\textbf{Inductive step $k \ge 3$.}  Assume Theorem~\ref{thm:main} for
$k - 1$ (in its rank-zero regime $n - 2 \ge 2(k-1)$, i.e., $n \ge 2k$):
$g_{k-1}(q) = (q+1)^{\binom{k-1}{2}}$.

By Lemma~\ref{lem:Ck},
\[
   R'_n\, v_{T_{\mathrm{first}}^{(\lambda)}}
   \;=\; (q+1)^{k-1} \cdot \Phi_*(v_{T_{\mathrm{first}}^{(\mu_1)}}).
\]
By the $\Phi_*$-commutation of the May-13 r1 paper (\texttt{2026-05-13-r1-inductive-reduction.tex}),
the iteration $R'_{j_0(\lambda)} \cdots R'_{n-1}$ pulls through $\Phi_*$
into the outer factor $M$:
\[
   R'_{j_0(\lambda)}\cdots R'_{n-1}\, \Phi_*(v_{T_{\mathrm{first}}^{(\mu_1)}})
   \;=\; M\, \Phi_*\!\left(R'^{(\mu_1)}_{j_0(\mu_1)}\cdots R'^{(\mu_1)}_{n-2}\,
       v_{T_{\mathrm{first}}^{(\mu_1)}}\right)
   \;=\; M\, \Phi_*(w_{\mu_1}),
\]
where $w_{\mu_1} := B^{(\mu_1)}_{j_0(\mu_1)}\, v_{T_{\mathrm{first}}^{(\mu_1)}}
\in V_{\mu_1}$.  Composing:
\[
   B^{(\lambda)}_{j_0(\lambda)}\, v_{T_{\mathrm{first}}^{(\lambda)}}
   \;=\; (q+1)^{k-1} \cdot M\, \Phi_*(w_{\mu_1}).
\]

Expand $w_{\mu_1} = \sum_{T' \in \mathrm{SYT}(\mu_1)} c_{T'}\, v_{T'}$.
By the inductive hypothesis, $c_{T_{*}^{(\mu_1)}} = g_{k-1}(q) =
(q+1)^{\binom{k-1}{2}}$.

Apply $\Phi_*$ linearly, then $M$, then extract the coefficient at
$v_{T_{*}^{(\lambda)}}$:
\[
   g_k(q) \;=\; (q+1)^{k-1} \cdot \sum_{T' \in \mathrm{SYT}(\mu_1)}
       c_{T'}\, \langle v_{T_{*}^{(\lambda)}},\, M\,\Phi_*(v_{T'})\rangle.
\]
By Lemma B of the hook-dim-1 paper (Step 2's combinatorial argument:
$M$ does not move entries $\le n-k$, and only $T' = T_{*}^{(\mu_1)}$
has its column-$1$ small entries match $T_{*}^{(\lambda)}$), the sum
collapses to the $T' = T_{*}^{(\mu_1)}$ term.  By Lemma~\ref{lem:fk},
the surviving inner product equals $1$.  Therefore
\[
   g_k(q) \;=\; (q+1)^{k-1} \cdot c_{T_{*}^{(\mu_1)}} \cdot 1
   \;=\; (q+1)^{k-1} \cdot (q+1)^{\binom{k-1}{2}}
   \;=\; (q+1)^{\binom{k-1}{2} + (k-1)}
   \;=\; (q+1)^{\binom{k}{2}},
\]
using the Pascal-like identity $\binom{k-1}{2} + (k-1) = \binom{k}{2}$.
\end{proof}

\section{Computational verification}\label{sec:verify}

\paragraph{Scripts.}  All verification scripts use \texttt{verify\_seminormal.py}'s
asymmetric Murphy normalization ($b' = 1$).

\begin{itemize}[leftmargin=2em,topsep=0.2em,itemsep=0.1em]
\item \texttt{\textasciitilde/projects/scratch/prove-2026-05-13-tracer-formula/verify\_formula.py}:
   computes $g_k(q)$ symbolically and confirms
   $g_k(q) = (q+1)^{\binom{k}{2}}$ for all tested $(k, n)$ with diff $= 0$.
\item \texttt{verify\_Ck\_fk.py}: verifies Lemmas~\ref{lem:Ck}
   and~\ref{lem:fk} directly, by computing the coefficient of $v_{T'_-}$
   in $R'_n v_{T_{\mathrm{first}}^{(\lambda)}}$ (expected
   $(q+1)^{k-1}$) and the coefficient of $v_{T_{*}^{(\lambda)}}$ in
   $M v_{T'_-}$ (expected $1$).
\end{itemize}

\paragraph{Verified cases (Theorem~\ref{thm:main}).}
$(k, n) \in \{(2, 4), (2, 5), (2, 6), (3, 6), (3, 7), (3, 8), (4, 8),
(4, 9), (5, 10)\}$.  Exact match of $g_k(q) = (q+1)^{\binom{k}{2}}$
in every case.

\paragraph{Verified cases (Lemmas~\ref{lem:Ck} and~\ref{lem:fk}).}
$(k, n) \in \{(2, 4), (2, 5), (3, 6), (3, 7), (3, 8), (4, 8), (4, 9),
(5, 10)\}$.  $C_k = (q+1)^{k-1}$ and $f_k = 1$ in every case.

\section{Honest status and consequences}\label{sec:status}

\paragraph{Proved unconditionally (in the asymmetric Murphy normalization
of the seminormal basis).}  $g_k(q) = (q+1)^{\binom{k}{2}}$ for every
hook $\lambda = (k, 1^{n-k})$ with $k \ge 2$, $n \ge 2k$.  This closes
Corollary~5.2 of \texttt{2026-05-13-hook-dim-equals-one.tex}, which had
verified the formula computationally up to $k = 5$ but left the proof
open.

\paragraph{Normalization-dependence.}  The exact constant $(q+1)^{\binom{k}{2}}$
depends on the choice of seminormal off-diagonals.  In the symmetric
Murphy convention ($b = b' = \sqrt{b_d b'_d}$ with consistent sign), the
formula picks up a product of $\sqrt{b_d b'_d}$ factors.  However, the
\emph{shape} of the formula --- a single monomial in $(q+1)$ with
exponent $\binom{k}{2}$ --- is structural.  Any seminormal normalization
yields $g_k(q) = (q+1)^{\binom{k}{2}} \cdot \prod (\text{Hoefsmit prefactors})$.

\paragraph{Why the answer is $n$-independent.}  The crucial structural
fact is that:
\begin{enumerate}[leftmargin=2em,topsep=0.2em,itemsep=0.1em,label=(\roman*)]
\item The same-row prefix in Lemma~\ref{lem:Ck} contributes $(q+1)^{k-1}$,
   independent of $n$ (the prefix has length $k - 1$, the number of
   same-row pairs in row 1 of $T_{\mathrm{first}}$).
\item The 2-block phase in Lemma~\ref{lem:Ck} ``slides'' the new-SYT
   coefficient unchanged across $n - k$ steps via the off-diagonal
   $b' = 1$, contributing no $n$-dependent factor.
\item The all-swap chain in Lemma~\ref{lem:fk} has length $k - 2$,
   independent of $n$.
\end{enumerate}
The $n$-dependent Hoefsmit prefactors $(a_d + 1)$ for $d = -k, -(k+1),
\ldots, -(n-1)$ all appear in the orthogonal $\alpha$-coefficients
$\alpha_j$ but NOT in the persistent $\beta_j = (q+1)^{k-1}$, because
the asymmetric Murphy off-diagonal $b'_d = 1$ is itself $d$-independent.
The structural ``rigidity'' of the formula is a feature of this
normalization choice.

\paragraph{Beyond hooks.}  For non-hook shapes with multiple corners,
$\dim B$ can exceed $1$, and the tracer-pair framework requires
generalization.  See the May-13 r1 paper for the iteration scheme on
$r = 1$ shapes; tracer pairs adapted to multi-corner shapes are
open.  The technique of refining ``$\ne 0$'' to an exact $(q+1)$-monomial
via $b' = 1$ propagation appears robust and may transfer.

\paragraph{Combinatorial reading.}  The exponent $\binom{k}{2}$ is
suggestive: it is the number of inversions in the longest word of
$S_k$.  The all-swap chain in Lemma~\ref{lem:fk} executes a
$(k-1)$-cycle of length $k - 2$, contributing $k - 2$ inversions to
the permutation $T'_- \mapsto T_{*}^{(\lambda)}$.  The $R'_n$ trace in
Lemma~\ref{lem:Ck} contributes $k - 1$ same-row pairs.  Summing
$(k-2) + (k-1) + \binom{k-1}{2} = \binom{k}{2}$ recovers the answer
inductively.  A direct combinatorial interpretation of the exponent
remains open.

\end{document}
